% =============================================================================
\section{Fuerzas electrostáticas sobre conductores}\label{sec:fuerzas}
% =============================================================================

Dos conductores cargados se ejercen fuerzas. En principio podrían calcularse con la ley de Coulomb, integrando la fuerza sobre cada
elemento de superficie, o con la presión electrostática $\sigma^2/2\eo$ del Tema~3; pero la distribución de carga cambia al mover los
conductores y el cálculo directo es casi siempre intratable. La alternativa es la conservación de la energía.

\subsection{El teorema de Earnshaw}

Si un conductor está en reposo y sobre él actúa una fuerza electrostática neta no nula, algo tiene que equilibrarla. El siguiente
teorema muestra que ese algo no puede ser otra fuerza electrostática.

\begin{teorema}[Earnshaw, 1842]\label{teo:earnshaw}
Una carga puntual no puede estar en equilibrio estable bajo la acción exclusiva de fuerzas electrostáticas~\cite{Earnshaw}. En consecuencia, ningún
sistema de cargas puede mantenerse en equilibrio estable solo por interacciones coulombianas.
\end{teorema}

\begin{proof}
Sea $q>0$ una carga situada en $\vr_0$, en una región sin otras cargas, y sea $\phi$ el potencial creado por las demás. La energía potencial de
$q$ es $q\phi(\vr)$. Para que $\vr_0$ fuese una posición de equilibrio estable, $q\phi$ debería tener allí un mínimo local estricto: cualquier
desplazamiento debería aumentar la energía, de modo que la fuerza devolviera la carga a $\vr_0$. Ahora bien, en esa región $\phi$ satisface la
ecuación de Laplace, $\lap\phi=0$, porque las cargas que lo crean están en otro sitio. Si $\phi$ tuviese un mínimo estricto en $\vr_0$, en sus
proximidades sería
\[
  \frac{\partial^2\phi}{\partial x^2}\geq0,\qquad\frac{\partial^2\phi}{\partial y^2}\geq0,\qquad\frac{\partial^2\phi}{\partial z^2}\geq0,
\]
y como su suma es nula, las tres deberían anularse. Más rigurosamente, por la propiedad del valor medio de las funciones armónicas~\cite{T3,Kellogg},
$\phi(\vr_0)$ es el promedio de $\phi$ sobre cualquier esfera centrada en $\vr_0$ contenida en la región, y un promedio no puede ser estrictamente
menor que todos los valores promediados: $\phi$ no tiene mínimos ni máximos locales estrictos donde no hay carga. Para $q<0$ se razona igual con los
máximos. Como el razonamiento vale para cada carga del sistema, el sistema tampoco puede estar en equilibrio estable.
\end{proof}

\begin{observacion}[Puntos de silla]\label{obs:silla}
Sí puede haber puntos de \emph{equilibrio}, donde $\vE=\vcero$; pero son puntos de silla del potencial: estables en unas direcciones e inestables
en otras (\cref{fig:earnshaw}). Por eso, junto a las fuerzas electrostáticas sobre un conductor en reposo, siempre debe actuar una fuerza
mecánica (sujeciones, muelles, la cohesión del propio material) que las equilibre. Ya lo vimos en el Tema~3: la presión electrostática tiende a
arrancar la carga de la superficie, y lo impide la ligadura del metal.
\end{observacion}

% Figura: teorema de Earnshaw. Potencial de cuatro cargas positivas iguales en los
% vértices de un cuadrado (z = 0, vértices (+-1, +-1)), en el plano y = 0 cerca del centro:
% mínimo a lo largo de x, máximo a lo largo de z  ->  punto de silla.
\begin{figure}[t]
\centering
\begin{subfigure}[b]{0.56\textwidth}\centering
\begin{tikzpicture}
\begin{axis}[width=7.6cm, height=6.0cm, view={-32}{24}, xmin=-0.8, xmax=0.8, ymin=-0.8, ymax=0.8,
   xlabel={$x/a$}, ylabel={$z/a$}, zlabel={$\phi/\phi_0$}, xtick={-0.5,0,0.5}, ytick={-0.5,0,0.5},
   ztick={2.6,2.8,3.0}, colormap={tema}{color=(acentoB) color=(acentoA) color=(acentoC)},
   z buffer=sort, clip=false, label style={font=\footnotesize},
   x label style={sloped}, y label style={sloped}]
  \addplot3[surf, shader=faceted interp, faceted color=fondo!70!texto, line width=0.1pt,
     samples=29, domain=-0.8:0.8, y domain=-0.8:0.8, opacity=0.95]
     {1/sqrt((x-1)^2+1+y^2)+1/sqrt((x+1)^2+1+y^2)+1/sqrt((x-1)^2+1+y^2)+1/sqrt((x+1)^2+1+y^2)};
  \addplot3[only marks, mark=*, mark size=1.6pt, color=texto] coordinates {(0,0,2.8284)};
\end{axis}
\end{tikzpicture}
\caption{$\phi$ en el plano $y=0$ cerca del centro}
\end{subfigure}\hfill
\begin{subfigure}[b]{0.42\textwidth}\centering
\begin{tikzpicture}
\begin{axis}[width=6.0cm, height=5.0cm, xmin=-0.8, xmax=0.8, ymin=2.45, ymax=3.15,
   xlabel={desplazamiento$/a$}, ylabel={$\phi/\phi_0$}, grid=major, ytick={2.5,2.75,3.0}]
  \addplot[curvaE, domain=-0.8:0.8, samples=101]
     {2/sqrt((x-1)^2+1)+2/sqrt((x+1)^2+1)};
  \addplot[curvaD, domain=-0.8:0.8, samples=101] {4/sqrt(2+x^2)};
  \node[etiqf, font=\scriptsize, text=campoE] at (axis cs:0,3.02) {según $x$: mínimo};
  \node[etiqf, font=\scriptsize, text=campoD] at (axis cs:0,2.6) {según $z$: máximo};
  \addplot[only marks, mark=*, mark size=1.5pt, color=texto] coordinates {(0,2.8284)};
\end{axis}
\end{tikzpicture}
\caption{Cortes por el centro}
\end{subfigure}
\caption{Teorema de Earnshaw. Cuatro cargas positivas iguales en los vértices de un cuadrado de lado $2a$ en el
plano $z=0$, con $\phi_0=q/(4\pi\eo a)$. En el centro del cuadrado el campo es nulo, pero el potencial no tiene allí un mínimo:
crece al desplazarse según $x$ (hacia dos de las cargas) y decrece según $z$ (al alejarse del plano). Es un punto de
silla, como exige $\lap\phi=0$: una carga positiva colocada en el centro está en equilibrio, pero inestable frente a
desplazamientos perpendiculares al plano.}
\label{fig:earnshaw}
\end{figure}


\subsection{El método de los trabajos virtuales}

Imaginemos que uno de los conductores sufre un desplazamiento infinitesimal $\dl$ (o un giro $\dd\theta$), lento y reversible, mientras los
demás permanecen fijos (\cref{fig:desplazamiento}). La fuerza electrostática sobre él realiza el trabajo
\begin{equation}\label{eq:trabajo-mec}
  \dd\Wmec=\vFe\cdot\dl \qquad(\text{o } \dd\Wmec=\vtau_e\cdot\dd\vec\theta \text{ en un giro}),
\end{equation}
y el balance de energía depende de qué se mantiene constante durante el desplazamiento. Estudiamos los dos casos más sencillos.

% Figura: método de los trabajos virtuales. Desplazamiento infinitesimal de un
% conductor (a) aislado (cargas constantes) y (b) conectado a baterías (potenciales constantes).
\begin{figure}[t]
\centering
\tikzset{
  cond/.style={fill=conductor, draw=texto, line width=0.6pt},
  condv/.style={draw=texto, dashed, line width=0.5pt, dash pattern=on 2.5pt off 1.8pt},
  bateria/.pic={
    \draw[texto, line width=0.6pt] (-0.35,0) -- (-0.08,0) (0.08,0) -- (0.35,0);
    \draw[texto, line width=1.1pt] (-0.08,-0.22) -- (-0.08,0.22);
    \draw[texto, line width=0.6pt] (0.08,-0.13) -- (0.08,0.13);}
}
\begin{subfigure}[b]{0.47\textwidth}\centering
\begin{tikzpicture}[scale=0.88]
  \draw[cond] plot[smooth cycle, tension=0.9] coordinates {(-2.3,0.1) (-1.9,0.75) (-1.2,0.6) (-1.0,-0.1) (-1.6,-0.6) (-2.2,-0.45)};
  \node[font=\small, text=black] at (-1.65,0.05) {$Q_1$};
  \draw[cond] plot[smooth cycle, tension=0.9] coordinates {(0.6,0.25) (1.0,0.85) (1.6,0.7) (1.75,0.0) (1.3,-0.55) (0.75,-0.4)};
  \node[etiq, anchor=south] at (1.2,0.9) {$Q_2$};
  \draw[condv] plot[smooth cycle, tension=0.9] coordinates {(1.05,0.25) (1.45,0.85) (2.05,0.7) (2.2,0.0) (1.75,-0.55) (1.2,-0.4)};
  \draw[vecP] (1.75,0.0) -- (2.2,0.0);
  \node[etiq, text=campoP, anchor=west] at (2.25,0.0) {$\dl$};
  \node[etiqt, anchor=north] at (-0.3,-1.0) {cargas aisladas: $\dd Q_i=0$};
  \node[etiq, anchor=north] at (-0.3,-1.45) {$\dd W_{\mathrm{mec}}+\dd U_e=0$};
\end{tikzpicture}
\caption{Sistema aislado}
\end{subfigure}\hfill
\begin{subfigure}[b]{0.47\textwidth}\centering
\begin{tikzpicture}[scale=0.88]
  \draw[cond] plot[smooth cycle, tension=0.9] coordinates {(-2.3,0.1) (-1.9,0.75) (-1.2,0.6) (-1.0,-0.1) (-1.6,-0.6) (-2.2,-0.45)};
  \node[font=\small, text=black] at (-1.65,0.05) {$\phi_1$};
  \draw[cond] plot[smooth cycle, tension=0.9] coordinates {(0.6,0.25) (1.0,0.85) (1.6,0.7) (1.75,0.0) (1.3,-0.55) (0.75,-0.4)};
  \node[etiq, anchor=south] at (1.2,0.9) {$\phi_2$};
  \draw[condv] plot[smooth cycle, tension=0.9] coordinates {(1.05,0.25) (1.45,0.85) (2.05,0.7) (2.2,0.0) (1.75,-0.55) (1.2,-0.4)};
  \draw[vecP] (1.75,0.0) -- (2.2,0.0);
  % cables y baterías
  \draw[texto, line width=0.55pt] (-1.65,-0.55) -- (-1.65,-1.0) -- (-0.75,-1.0);
  \pic at (-0.4,-1.0) {bateria};
  \draw[texto, line width=0.55pt] (-0.05,-1.0) -- (0.25,-1.0);
  \draw[texto, line width=0.55pt] (1.15,-0.5) -- (1.15,-1.0) -- (0.55,-1.0);
  \pic at (0.4,-1.0) {bateria};
  \node[etiq, text=campoP, anchor=west] at (2.25,0.0) {$\dl$};
  \node[etiqt, anchor=north] at (-0.1,-1.35) {potenciales fijos: $\dd\phi_i=0$};
  \node[etiq, anchor=north] at (-0.1,-1.8) {$\dd W_{\mathrm{mec}}+\dd U_e=\dd W_{\mathrm{bat}}$};
\end{tikzpicture}
\caption{Sistema a potencial constante}
\end{subfigure}
\caption{Método de los trabajos virtuales. Se imagina un desplazamiento infinitesimal $\dl$ de uno de los conductores
(trazos) y se hace el balance de energía. (a) Si el sistema está aislado, las cargas no cambian y el trabajo mecánico
$\dd W_{\mathrm{mec}}=\vF_e\cdot\dl$ sale de la energía electrostática. (b) Si los conductores están unidos a baterías, sus potenciales
no cambian, pero las baterías deben suministrar las cargas $\dd Q_i$, con un trabajo $\dd W_{\mathrm{bat}}=\sum_i\phi_i\,\dd Q_i$.}
\label{fig:desplazamiento}
\end{figure}


\begin{teorema}[Fuerza a carga constante]\label{teo:fuerza-Q}
Si el sistema de conductores está aislado, de modo que las cargas $Q_i$ permanecen constantes, la fuerza electrostática y el momento sobre un
conductor son
\begin{equation}\label{eq:fuerza-Q}
  \boxed{\;\vFe=-\grad U_e\big|_{Q}\;},\qquad F_\xi=-\frac{\partial U_e}{\partial\xi}\bigg|_{Q},\qquad \tau_\theta=-\frac{\partial U_e}{\partial\theta}\bigg|_{Q},
\end{equation}
donde el gradiente se toma respecto a las coordenadas que describen la posición del conductor. En función de los coeficientes de potencial,
\begin{equation}\label{eq:fuerza-pij}
  F_\xi=-\frac12\sum_{i,j}Q_iQ_j\,\frac{\partial p_{ij}}{\partial\xi} .
\end{equation}
\end{teorema}

\begin{proof}
Al estar aislado, no entra energía de fuera: el trabajo mecánico que realiza la fuerza electrostática se obtiene a costa de la energía
almacenada,
\[
  \dd\Wmec+\dd U_e=0\quad\Longrightarrow\quad \vFe\cdot\dl=-\dd U_e\big|_Q .
\]
Como $U_e$ es función de la configuración, $\dd U_e=\grad U_e\cdot\dl$ y, al ser $\dl$ arbitrario, $\vFe=-\grad U_e|_Q$. Para un giro, $\tau_\theta\,\dd\theta=-\dd U_e|_Q$. Por
último, con $U_e=\frac12\sum p_{ij}Q_iQ_j$ y las $Q_i$ constantes, solo dependen de la posición los coeficientes $p_{ij}$, que son puramente geométricos, y se obtiene
\eqref{eq:fuerza-pij}.
\end{proof}

\begin{teorema}[Fuerza a potencial constante]\label{teo:fuerza-phi}
Si cada conductor está unido a una batería que mantiene fijo su potencial $\phi_i$, la fuerza y el momento sobre un conductor son
\begin{equation}\label{eq:fuerza-phi}
  \boxed{\;\vFe=+\grad U_e\big|_{\phi}\;},\qquad F_\xi=+\frac{\partial U_e}{\partial\xi}\bigg|_{\phi},\qquad \tau_\theta=+\frac{\partial U_e}{\partial\theta}\bigg|_{\phi},
  \qquad F_\xi=\frac12\sum_{i,j}\phi_i\phi_j\,\frac{\partial c_{ij}}{\partial\xi}.
\end{equation}
\end{teorema}

\begin{proof}
Ahora el sistema no está aislado. Para mantener los potenciales al desplazar el conductor, las baterías deben suministrar cargas $\dd Q_i$, trayéndolas
desde un potencial nulo hasta $\phi_i$; por el \cref{lem:trabajo}, su trabajo es
\[
  \dd\Wbat=\sum_i\phi_i\,\dd Q_i .
\]
La variación de energía electrostática, con los $\phi_i$ fijos, es $\dd U_e=\dd\big(\frac12\sum_i\phi_iQ_i\big)=\frac12\sum_i\phi_i\,\dd Q_i=\frac12\dd\Wbat$. El balance de energía es
\[
  \begin{aligned}
  \dd\Wmec+\dd U_e=\dd\Wbat\quad\Longrightarrow\quad \dd\Wmec&=\dd\Wbat-\dd U_e=\sum_i\phi_i\,\dd Q_i-\frac12\sum_i\phi_i\,\dd Q_i\\
  &=\frac12\sum_i\phi_i\,\dd Q_i=\dd U_e\big|_\phi .
  \end{aligned}
\]
Así $\vFe\cdot\dl=+\dd U_e|_\phi$ y $\vFe=+\grad U_e|_\phi$. La última forma se obtiene de $U_e=\frac12\sum c_{ij}\phi_i\phi_j$ con los $\phi_i$ constantes.
\end{proof}

\begin{observacion}[El reparto de la energía de las baterías]\label{obs:baterias}
A potencial constante, de la energía que aportan las baterías, la mitad se convierte en trabajo mecánico y la otra mitad aumenta la energía
almacenada en el campo. Por eso el signo cambia: el sistema evoluciona hacia estados de \emph{mayor} energía electrostática, porque las baterías
pagan el doble de lo que se almacena. A carga constante, en cambio, el sistema tiende al mínimo de $U_e$.
\end{observacion}

\subsection{Las dos fórmulas dan la misma fuerza}

El signo opuesto de \eqref{eq:fuerza-Q} y \eqref{eq:fuerza-phi} podría hacer pensar que la fuerza depende del proceso. No es así: la fuerza sobre un
conductor depende solo de la configuración en un instante, de las cargas y de los campos que hay en ese momento, no de lo que vaya a ocurrir en el
desplazamiento imaginario.

\begin{proposicion}[Equivalencia de ambos métodos]\label{prop:equivalencia}
En una misma configuración, con cargas $Q_i$ y potenciales $\phi_i$ relacionados por $\boldsymbol\phi=\mathsf P\boldsymbol Q$,
\begin{equation}\label{eq:equivalencia}
  -\frac{\partial U_e}{\partial\xi}\bigg|_{Q}=+\frac{\partial U_e}{\partial\xi}\bigg|_{\phi} .
\end{equation}
\end{proposicion}

\begin{proof}
Sea $\mathsf P(\xi)$ la matriz de coeficientes de potencial y $\mathsf C(\xi)=\mathsf P(\xi)^{-1}$ la de capacidad. Derivando $\mathsf C\mathsf P=\mathsf 1$ respecto a $\xi$ se obtiene
$\mathsf C'\mathsf P+\mathsf C\mathsf P'=\mathsf 0$, es decir, $\mathsf C'=-\mathsf C\mathsf P'\mathsf C$. Entonces, con $\boldsymbol Q=\mathsf C\boldsymbol\phi$ y la simetría de $\mathsf C$,
\[
  \frac{\partial U_e}{\partial\xi}\bigg|_{\phi}=\frac12\,\boldsymbol\phi^{\mathsf T}\mathsf C'\boldsymbol\phi=-\frac12\,\boldsymbol\phi^{\mathsf T}\mathsf C\mathsf P'\mathsf C\boldsymbol\phi
  =-\frac12\,\boldsymbol Q^{\mathsf T}\mathsf P'\boldsymbol Q=-\frac{\partial U_e}{\partial\xi}\bigg|_{Q}. \qedhere
\]
\end{proof}

\begin{corolario}[Fuerza sobre una armadura de condensador]\label{cor:fuerza-condensador}
Si la capacidad de un condensador depende de una coordenada $\xi$, la fuerza generalizada sobre la armadura móvil es
\begin{equation}\label{eq:fuerza-C}
  \boxed{\;F_\xi=\frac12\,V^2\,\frac{\dd C}{\dd\xi}=\frac{Q^2}{2C^2}\,\frac{\dd C}{\dd\xi}\;},
\end{equation}
y el momento, $\tau_\theta=\frac12V^2\,\dd C/\dd\theta$. El sistema tiende a aumentar su capacidad.
\end{corolario}

\begin{proof}
A carga constante, $F_\xi=-\dfrac{\dd}{\dd\xi}\Big(\dfrac{Q^2}{2C}\Big)=\dfrac{Q^2}{2C^2}\dfrac{\dd C}{\dd\xi}$; a potencial constante, $F_\xi=\dfrac{\dd}{\dd\xi}\Big(\dfrac12CV^2\Big)=\dfrac12V^2\dfrac{\dd C}{\dd\xi}$.
Ambas coinciden porque $Q=CV$.
\end{proof}

\begin{metodo}{Receta: fuerza por trabajos virtuales}
\begin{pasos}
  \item Elegir la coordenada $\xi$ que describe el movimiento posible (distancia, longitud introducida, ángulo).
  \item Expresar la energía en función de $\xi$: $U_e=Q^2/2C(\xi)$ o $\frac12C(\xi)V^2$ para un condensador; en general $\frac{\eo}{2}\int E^2\dd V$ o
        $\frac12\sum p_{ij}Q_iQ_j$.
  \item Derivar con el signo que corresponda: $-\partial_\xi U_e$ a carga constante, $+\partial_\xi U_e$ a potencial constante.
  \item Interpretar el signo: $F_\xi>0$ significa que la fuerza tiende a aumentar $\xi$.
\end{pasos}
\end{metodo}

\begin{ejemplo}[Comprobación con la presión electrostática]\label{ej:presion}
En un condensador plano de área $S$ con cargas $\pm Q$ y separación $\xi$, $C=\eo S/\xi$ y $F_\xi=\frac{Q^2}{2C^2}\frac{\dd C}{\dd\xi}=-\frac{Q^2}{2\eo S}$: las placas se atraen.
Con la presión del Tema~3, cada placa sufre $P=\sigma^2/2\eo$ dirigida hacia el exterior del metal, es decir, hacia la otra placa:
$F=\frac{\sigma^2}{2\eo}S=\frac{Q^2}{2\eo S}$. Coinciden (\cref{prob:fuerza-placas}).
\end{ejemplo}

\begin{ejemplo}[El acelerómetro capacitivo]\label{ej:acelerometro}
Los acelerómetros de los teléfonos móviles son condensadores diferenciales microfabricados: una placa unida a una masa sísmica se mueve
entre dos placas fijas, y la aceleración se mide por el desequilibrio $C_1-C_2$ (\cref{fig:acelerometro}). Las fuerzas electrostáticas sobre la
placa móvil se calculan exactamente como en el \cref{prob:tres-placas}; y, por el teorema de Earnshaw, son los muelles los que proporcionan un
equilibrio estable.
\end{ejemplo}

% Figura: acelerómetro capacitivo (condensador diferencial). Masa sísmica móvil
% entre dos placas fijas; capacidades C1 = eps0 S/(d - x), C2 = eps0 S/(d + x).
\begin{figure}[t]
\centering
\begin{subfigure}[b]{0.5\textwidth}\centering
\begin{tikzpicture}[scale=1.0]
  \tikzset{muelle/.style={decorate, decoration={coil, aspect=0.45, segment length=2.4mm, amplitude=1.6mm},
                          draw=texto!80!fondo, line width=0.55pt}}
  % bastidor
  \fill[conductor!55!fondo] (-2.9,-1.35) rectangle (-2.7,1.35);
  \fill[conductor!55!fondo] (2.7,-1.35) rectangle (2.9,1.35);
  \draw[texto, line width=0.5pt] (-2.9,-1.35) rectangle (-2.7,1.35) (2.7,-1.35) rectangle (2.9,1.35);
  % masa sísmica (placa central móvil) y muelles
  \fill[conductor] (-0.45,-0.55) rectangle (0.45,0.55);
  \draw[texto, line width=0.6pt] (-0.45,-0.55) rectangle (0.45,0.55);
  \node[font=\small, text=black] at (0,0) {$m$};
  \draw[muelle] (-2.7,0) -- (-0.45,0);
  \draw[muelle] (0.45,0) -- (2.7,0);
  % placas fijas del condensador diferencial (debajo)
  \draw[placa] (-1.0,-1.25) -- (-1.0,-2.25);
  \draw[placa] (1.0,-1.25) -- (1.0,-2.25);
  \draw[placa, draw=acentoC] (0,-0.55) -- (0,-2.25);
  \node[etiqt, anchor=north] at (-1.0,-2.3) {fija};
  \node[etiqt, anchor=north] at (1.0,-2.3) {fija};
  \node[etiq, text=acentoC, anchor=north] at (0,-2.3) {móvil};
  \node[etiq] at (-0.5,-1.5) {$C_1$};
  \node[etiq] at (0.5,-1.5) {$C_2$};
  \draw[cota] (-1.0,-1.95) -- (0,-1.95) node[midway, below, etiq, font=\footnotesize] {$d$};
  \draw[cota] (0,-1.95) -- (1.0,-1.95) node[midway, below, etiq, font=\footnotesize] {$d$};
  % aceleración
  \draw[vecE] (-0.6,1.0) -- (0.6,1.0);
  \node[etiq, text=campoE, anchor=south] at (0,1.05) {$\vec a$};
\end{tikzpicture}
\caption{Esquema}
\end{subfigure}\hfill
\begin{subfigure}[b]{0.46\textwidth}\centering
\begin{tikzpicture}
\begin{axis}[width=6.4cm, height=5.0cm, xmin=-0.6, xmax=0.6, ymin=-2.0, ymax=2.6, grid=major,
   xlabel={$x/d$}, ylabel={$C\,d/(\eo S)$}, xtick={-0.5,0,0.5}, ytick={-2,-1,0,1,2},
   legend style={at={(0.5,1.03)}, anchor=south, legend columns=3, font=\scriptsize}]
  \addplot[curvaE, domain=-0.6:0.6, samples=80] {1/(1-x)}; \addlegendentry{$C_1$}
  \addplot[curvaD, domain=-0.6:0.6, samples=80] {1/(1+x)}; \addlegendentry{$C_2$}
  \addplot[curvaP, domain=-0.6:0.6, samples=80] {1/(1-x)-1/(1+x)}; \addlegendentry{$C_1-C_2$}
\end{axis}
\end{tikzpicture}
\caption{Capacidades frente al desplazamiento}
\end{subfigure}
\caption{Acelerómetro capacitivo (sistema MEMS). Una masa sísmica sujeta por muelles lleva una placa que se
mueve entre dos placas fijas; en reposo, a distancia $d$ de ambas. Una aceleración la desplaza $x$ y desequilibra el
condensador diferencial: $C_1-C_2=\frac{\eo S}{d-x}-\frac{\eo S}{d+x}\simeq\frac{2\eo S}{d^2}\,x$ para $x\ll d$. Las fuerzas electrostáticas
entre las placas, que se calculan como en el \cref{prob:tres-placas}, deben ser pequeñas frente a las de los muelles,
porque el equilibrio electrostático puro es inestable (\cref{teo:earnshaw}).}
\label{fig:acelerometro}
\end{figure}

